\subsection{Définition structurelle et fonction des constantes de Planck}
\label{sec:ii-definicion-planck}

\begin{definition}[Détermination structurelle de l'action angulaire]
Dans la droite orientée \(\mathcal L_S\), la détermination HMT de
l'action angulaire est le couple de sections
\(\hbar_{\mathrm{pre}},\hbar_{\mathrm{ret}}\) de
\eqref{eq:el-secciones-accion}, avec leur identification par retour.
Sa structure réunit la valuation de la norme locale
\(\varphi\mathrm N_{G_H}(\alpha)\), le caractère \(H_5\) et la
continuation régionale \(\mathcal C_\pi\). L'action par tour complet
est \(h_a=2\pi\hbar_a\) dans chaque orientation \(a\).
\end{definition}

Cette définition conserve trois déterminations successives. Le quotient
local de \eqref{eq:el-h4} fixe les seize feuillets de la norme ;
l'inégalité \eqref{eq:el-decada} fixe son ordre décimal. Le caractère et
le renouvellement normalisé de \eqref{eq:revision-renovacion} déterminent
ensuite les coordonnées avant et après le retour. L'unicité de la queue
utilise son premier poids et sa règle de concaténation ; la distribution des
coefficients de \(H_5\) conserve le caractère déclaré à la
section~\ref{sec:revision-planck}. Dans la branche positive, l'information
multiplicative du retour est l'invariant
\begin{equation}
 \frac{\hbar_{\mathrm{pre}}}{\hbar_{\mathrm{ret}}}
 =R_{\mathrm{act}}=\frac{H_5}{\eta_{\mathrm{ret}}}>1.
 \label{eq:ii-definicion-planck-retorno}
\end{equation}

Sa fonction physique s’exprime au moyen de la contraction entre action et
fréquence : une section \(\hbar_a\) assigne une énergie à une fréquence
angulaire, tandis que \(h_a\) exprime l’action d’un tour. Le quotient
\eqref{eq:ii-definicion-planck-retorno} conserve la modification
structurelle entre les deux lectures même lorsque l’on change la base commune
d’unités. L’identification à la constante physique et sa précision
sont examinées dans la coordinatisation de la section~\ref{sec:revision-planck-contraste},
avec les deux coordonnées de retour différenciées.
